\documentclass[11pt]{article}
\usepackage[a4paper,margin=27mm]{geometry}
\usepackage{amsmath,amssymb,amsthm,mathtools}
\usepackage[hidelinks]{hyperref}
\newtheorem{theorem}{Theorem}
\newtheorem{lemma}[theorem]{Lemma}
\DeclareMathOperator{\Tr}{Tr}
\DeclareMathOperator{\ran}{ran}
\newcommand{\HS}{\mathrm{HS}}
\title{A trace inequality for four canonically normalized projections}
\author{Research proof and Lean formalization}
\date{October 27, 2026}
\begin{document}
\maketitle
\begin{abstract}
We prove the four-projection inequality stated as Conjecture 6.6 in
Banica--Nechita, arXiv:1602.04456v3. The canonical normalization uses the positive
Moore--Penrose inverse square root of the sum of each pair of projections.
The proof applies to arbitrary finite dimensions and does not require either
of the two rank-balance assumptions in that conjecture. Its main step is an
entrywise estimate in two independent spectral bases, followed by row and
column bounds for rectangular contractions.
\end{abstract}

\section{Statement and scope}
For a positive semidefinite matrix $T$, write $T^{\dagger/2}$ for the
positive operator with eigenvalue $t^{-1/2}$ at every positive eigenvalue $t$
of $T$, and eigenvalue zero on its kernel. All traces below are ordinary,
unnormalized traces. The displayed traces of products of orthogonal
projections are real and nonnegative.

\begin{theorem}\label{thm:main}
Let $P,Q,R,S$ be orthogonal projections on a finite-dimensional complex
Hilbert space. Suppose
\[
 PQ=RS=0,\qquad
 \ran P\cap\ran R=\ran Q\cap\ran S=\{0\}.
\]
Put $D=(P+R)^{\dagger/2}$ and $E=(Q+S)^{\dagger/2}$, and define
\[
 P'=DPD,\quad R'=DRD,\quad Q'=EQE,\quad S'=ESE.
\]
Then these four normalized operators are orthogonal projections, and
\begin{equation}\label{eq:target}
 \Tr(P'Q')+\Tr(R'S')\leq \Tr(PR)+\Tr(QS).
\end{equation}
\end{theorem}

Conjecture 6.6 of \cite[p.~20]{BN} additionally assumes
$\operatorname{rank}P+\operatorname{rank}Q=
\operatorname{rank}R+\operatorname{rank}S$ and
$\operatorname{rank}P+\operatorname{rank}R=
\operatorname{rank}Q+\operatorname{rank}S$.
Neither equality is needed here. The source proves the case $S=0$ in its
Proposition 6.7. We make no claim about the general Sinkhorn-convergence
questions elsewhere in that paper, strictness, or equality cases.

\section{The scalar estimate}
\begin{lemma}\label{lem:scalar}
Let $-1<a,b<1$, and set
$u=\sqrt{(1+a)(1+b)}$, $v=\sqrt{(1-a)(1-b)}$.
For complex numbers $k,l$ satisfying $uk+vl=0$,
\begin{equation}\label{eq:scalar}
 |k+l|^2\leq(a^2+b^2)(|k|^2+|l|^2).
\end{equation}
\end{lemma}
\begin{proof}
Both $u$ and $v$ are positive. Substitution of $l=-uk/v$ gives
\[
 |k+l|^2=
 \left(1-\frac{\sqrt{(1-a^2)(1-b^2)}}{1+ab}\right)
 (|k|^2+|l|^2).
\]
Write $t=a^2+b^2$ and $w=\sqrt{(1-a^2)(1-b^2)}$.
Since $1+ab>0$, it suffices to show $w\geq(1-t)(1+ab)$.
This is immediate if $t\geq1$. If $t<1$ and $ab<0$, then
$w^2=1-t+a^2b^2\geq1-t\geq(1-t)^2$, so
$w\geq1-t\geq(1-t)(1+ab)$. Finally, if $t<1$ and $ab\geq0$, the identity
\[
 (1-a^2)(1-b^2)-(1-a^2-b^2+ab)^2
  =(a-b)^2(1-a^2-b^2)
\]
implies $w\geq1-t+ab\geq(1-t)(1+ab)$.
\end{proof}

\section{Frames, grading, and normalization}
Choose orthonormal column frames $p,r,q,s$ for the four ranges, and write
\[
 X=[p\ r],\quad Y=[q\ s],\qquad G=X^*X,\quad H=Y^*Y.
\]
Both $X$ and $Y$ are injective. Indeed, $px+ry=0$ puts $px=-ry$ in the
intersection of the two ranges; disjointness and the isometry of each frame
give $x=y=0$. Thus $G,H$ are positive definite on their respective coordinate
spaces. Define
\[
 U=XG^{-1/2},\quad V=YH^{-1/2},\qquad K=U^*V.
\]
Here $U^*U=I$ and $V^*V=I$; no equality between their column counts is used.
In zero-dimensional coordinate spaces these statements have their usual
empty-matrix meanings.

The precise ambient normalization is obtained from
\begin{equation}\label{eq:polar}
 (XX^*)^{\dagger/2}X=X(X^*X)^{-1/2}=U.
\end{equation}
For completeness, a reduced singular-value decomposition $X=W\Sigma Z^*$,
with positive diagonal $\Sigma$, gives the two sides as $WZ^*$; the
ambient inverse square root is $W\Sigma^{-1}W^*$ and vanishes on
$\ran(W)^\perp$. The empty-column case is immediate.
Since $XX^*=P+R$, the columns of $U=[u_p\ u_r]$ are exactly
$u_p=Dp$ and $u_r=Dr$. Consequently
$P'=u_pu_p^*$ and $R'=u_ru_r^*$, with analogous formulas for $Q',S'$.
The columns of $U$ are orthonormal, so these are indeed orthogonal
projections. In the corresponding block decomposition of $K$,
\begin{equation}\label{eq:pinching}
 \Tr(P'Q')+\Tr(R'S')
 =\|K_{11}\|_{\HS}^2+\|K_{22}\|_{\HS}^2.
\end{equation}
This uses $\Tr(aa^*bb^*)=\|a^*b\|_{\HS}^2$.

Let $J=\operatorname{diag}(I,-I)$ on the $(p,r)$ coordinates, and let $L$
be its counterpart on the $(q,s)$ coordinates. Put $A=G-I$, $B=H-I$.
The isometry of each individual frame gives
\[
 A=\begin{pmatrix}0&p^*r\\r^*p&0\end{pmatrix},\qquad
 B=\begin{pmatrix}0&q^*s\\s^*q&0\end{pmatrix},
 \qquad JAJ=-A,\quad LBL=-B.
\]
In particular, $I-A=JGJ$ and $I-B=LHL$ are positive definite. The
eigenvalues of $A$ and $B$ therefore belong to $(-1,1)$.
The assumptions $PQ=RS=0$ give $p^*q=r^*s=0$, and hence
\[
 J(X^*Y)L=-X^*Y.
\]
Using $X=UG^{1/2}$, $Y=VH^{1/2}$ and the grading identities for the positive
square roots, we obtain, with $K'=JKL$,
\begin{equation}\label{eq:roots}
 (I+A)^{1/2}K(I+B)^{1/2}
 +(I-A)^{1/2}K'(I-B)^{1/2}=0.
\end{equation}

\section{Summation and conclusion}
Diagonalize $A$ and $B$ in independent orthonormal bases, with eigenvalues
$a_i$ and $b_j$. Write $k_{ij},l_{ij}$ for the entries of $K,K'$ in these
bases. Equation~\eqref{eq:roots} is exactly the hypothesis of
Lemma~\ref{lem:scalar} for each pair $(i,j)$. Therefore
\begin{equation}\label{eq:summed}
\begin{aligned}
 \tfrac14\|K+K'\|_{\HS}^2
 &\leq\tfrac14\sum_{i,j}(a_i^2+b_j^2)
                       (|k_{ij}|^2+|l_{ij}|^2)\\
 &\leq\tfrac12\left(\sum_i a_i^2+\sum_j b_j^2\right).
\end{aligned}
\end{equation}
For the last step, $K=U^*V$ is a rectangular contraction because $U,V$
are isometries; $K'=JKL$ is also a contraction. Thus every row and column
of each matrix has squared norm at most one, including after the two
independent unitary changes of basis. The last expression in
\eqref{eq:summed} is $(\Tr A^2+\Tr B^2)/2$.
The off-diagonal block descriptions above give
\[
 \tfrac12\Tr A^2=\|p^*r\|_{\HS}^2=\Tr(PR),\qquad
 \tfrac12\Tr B^2=\|q^*s\|_{\HS}^2=\Tr(QS).
\]
Finally, $(K+JKL)/2$ keeps precisely the two diagonal blocks of $K$.
Combining this observation with \eqref{eq:pinching} and
\eqref{eq:summed} proves Theorem~\ref{thm:main}.

\section{Certificate and research status}
The accompanying development uses Lean 4.34.1 and mathlib revision
\begin{center}
\texttt{d13f23b723b8a846827a245b89c10fc7d3f11612}.
\end{center}
The scalar estimate, independent spectral-basis comparison, genuine
continuous-functional-calculus inverse square roots, canonical range frames,
and complete frame-level inequality have been implemented. The final
projection-level assembly and independent full-certificate review are pending
for this manuscript revision; no completed source-level certificate is claimed
here. Computational checks are supplementary and do not replace the proof.

The source statement used here is specifically arXiv version 3 of \cite{BN}.
Comparison with the journal body remains incomplete. The documented bounded
literature search has found no prior full resolution, but includes inaccessible
material and does not establish worldwide priority. Polar decomposition,
spectral calculus, and contraction estimates are standard tools. This draft
remains unaccepted in the discovery campaign pending its stated gates.

\begin{thebibliography}{1}
\bibitem{BN} T. Banica and I. Nechita,
\emph{Flat matrix models for quantum permutation groups},
arXiv:1602.04456v3 (2016), Conjecture 6.6 and Proposition 6.7.
Journal reference: Adv. Appl. Math. \textbf{83} (2017), 24--46,
\href{https://doi.org/10.1016/j.aam.2016.09.001}{doi:10.1016/j.aam.2016.09.001}.
\end{thebibliography}
\end{document}
